\documentclass[12pt, a4paper]{article}
\usepackage{graphicx}
\usepackage{authblk}
\usepackage[margin=2.5cm]{geometry}
\usepackage[slovene]{babel}
\usepackage{hyperref}
\usepackage{listings}

\renewcommand{\contentsname}{Kazalo}
\addto\captionslovene{\renewcommand{\figurename}{Slika}}

\title{Linux spletni strežnik}
\author{REDACTED}
\date{Februar 2023}
\affil{Vegova Ljubljana}

\begin{document}

\begin{titlepage}
	\begin{center}
		\includegraphics[width=\textwidth]{slike/logotip_vegova_leze_brezokvirja.png}\par\vspace{6cm}
		{\scshape\huge Linux spletni strežnik}\par
		\large{Seminarska naloga pri informatiki}
		\vfill
		Mentor: Aleksander Lazarević \hfill Avtor: Maj Srčnik\par\vspace{1cm}
		Ljubljana, Februar 2023
	\end{center}
\end{titlepage}
\thispagestyle{empty}
\pagebreak
\tableofcontents
\pagebreak
\thispagestyle{empty}
\section*{Povzetek}
\section*{Abstract}
\pagebreak
\section{Uvod}
V modernih korporativnih medmrežjih je internetna samostojnost redek in pozabljen pojav, zato me neodvisne spletne strani vedno znova razveselijo.
V upanju da bo mogoče kdaj tudi moj skromni doprinos koga kratkočasil, sem sem se odločil postaviti lastni, neodvisni spletni strežnik.

Gostiti strežnik v lastnem domu je skladno z mojimi ideološkimi standardi, hkrati pa je odlična priložnost za poglobitev znanja na področju, ki je v učnem programu le površno obravnavan.

\section{Strojna oprema}
Ideja za lasten strežnik se mi je utrnila, ko mi je bil podarjen zdaj že 11 let star Raspberry Pi model B mikroračunalnik.
Na tej za dlan veliki ploščici so, od specifikacij ki me znaimajo, 700MHz enojedrni ARM procesor, 512GB RAM pomnilnika in 2 USB 2.0 priključka.
Razvidno je da Raspberry Pi ni konkurenten celo zelo ugodnim račualnikom istega namenskega razreda.
Kljub temu nisem imel na izbiro nič drugega in z zelo majno predvideno porabo (3.5W) je za trenutno energetsko krizo kar primeren.

Sistem je naložen na majhno microSD kartico, za katero upam da jo neprestani cikli branja in pisanja ne pokončajo, vsaj do konca pisanja seminarske naloge.
\begin{figure}[h]
	\centering
	\includegraphics[width=0.5\textwidth]{slike/raspberry.jpg}
	\caption{Raspberry Pi model B}
	\label{fig:raspberry}
\end{figure}

\section{Programska oprema}
\subsection{Operacijski sistem}
Raspbian OS je GNU/Linux operacijski sistem. Izbral sem ga, ker je zasnova na Debianu in doumevno optimiziran za Raspberry Pi računalnike.
GNU/Linux je zanesljiv, robusten in ekstensibilen.
Ta distribucija pa naj bi bila še posebaj nezahtevna za računska in pomnilniška sredstva\cite{raspbian}, zato je ustrezna izbira za moje projekt.

Ker strežnik ne bo nikoli priključen v ekran, za grafični namestitveni program sistema (installer), pa je ta vseeno potreben, sem ga začasno priklopil v hdmi priključek na televiziji in tako izvedel namestitev.

Debian vsebuje apt upravitelj paketov, zato nalaganje in nameščanje potrebnih programov ne bo velik problem.
V pomoč sta prav tako že namščena skrbnik prikritih procesov (daemon) systemd in strežnik za oddlajen dostop po protokolu SSH. 

\subsection{Spletni strežnik}
Po prijavi sem posodobil strežnik z \verb|sudo apt update && sudo apt upgrade| in namestil Nginx z \verb|sudo apt install nginx|.
Nginx je zmogljiv http strežnik, ki sicer ponuja še mnogo več, a jaz ga bom uporabljal le za http strežnik.\cite{nginx}

Ker nevem če bo na istem strežniku v prihodnsoti še kakšna stran, pišem konfiguracijo v \verb|/etc/nginx/sites-available/srcnik|.
Ker bi celotna vsebina nginx konfiguracijske datoteke lahko razkrila ranljivosti, bom navedel le del, ki je odgovoren za osebno spletno stran.
\begin{lstlisting}
	server{
		server_name srcnik.si www.srcnik.si
		location / {
			root /var/www/srcnik;
		}
	}
\end{lstlisting}

\section{Zaključek}
\bibliography{main}
\bibliographystyle{plain}
\section*{Izjava o avtorstvu}

\end{document}
